% ---
% school: "college"
% level: "8eme"
% section: null
% chapter: "العمليات والحساب على الأعداد الكسرية"
% source: "Axiomaths"
% title: "تحدّي : جذاء كسور"
% country: "Tunisie"
% year: 2026
% professor: "Axiomaths"
% difficulty: "Difficile"
% points: null
% tags: ["defi", "fractions", "produit-telescopique", "encadrement"]
% ---
% Exercice 18/18 de la série 8ème « العمليات والحساب على الأعداد الكسرية » (statement + \axsol + solution). Compiler avec XeLaTeX + axiomath.sty (+ axiomaths-logo.png).
\documentclass[12pt,a4paper]{article}
\usepackage[level=8b,color=15803D]{axiomath}
\axsetup{title={العمليات والحساب على الأعداد الكسرية}, corrige}

\begin{document}

\begin{corr}[title={تحدّي : جذاء كسور}, stars=4, source={Axiomaths}]
نعتبر العددين :
\[A=\dfrac{1}{2}\times\dfrac{3}{4}\times\dfrac{5}{6}\times\cdots\times\dfrac{99}{100}\qquad\qquad B=\dfrac{2}{3}\times\dfrac{4}{5}\times\dfrac{6}{7}\times\cdots\times\dfrac{100}{101}\]
(في كل منهما خمسون عاملا).
\begin{enumerate}
  \item قارن $\dfrac{1}{2}$ و $\dfrac{2}{3}$ ، ثم $\dfrac{3}{4}$ و $\dfrac{4}{5}$. بين أنه من أجل كل عدد طبيعي فردي $n$ :
  \[\dfrac{n}{n+1}<\dfrac{n+1}{n+2}\]
  \item استنتج أن $A<B$.
  \item احسب $A\times B$.
  \item استنتج أن $A\times A<\dfrac{1}{101}$ ، ثم أن $A<\dfrac{1}{10}$.
\end{enumerate}
\axsol
\begin{enumerate}[itemsep=14pt]
  \item نوحّد المقامات :
  \[\dfrac{1}{2}=\dfrac{3}{6}<\dfrac{4}{6}=\dfrac{2}{3}\qquad\qquad \dfrac{3}{4}=\dfrac{15}{20}<\dfrac{16}{20}=\dfrac{4}{5}\]
  وفي الحالة العامة، نكتب كل كسر على شكل $1$ ناقص كسر :
  \[\dfrac{n}{n+1}=1-\dfrac{1}{n+1}\qquad\qquad \dfrac{n+1}{n+2}=1-\dfrac{1}{n+2}\]
  \begin{prop}[مقلوبا عددين]
  إذا كان $0<u<v$ فإن $\dfrac{1}{u}>\dfrac{1}{v}$.
  \end{prop}
  $n+1<n+2$ ، إذن $\dfrac{1}{n+1}>\dfrac{1}{n+2}$ ، ومنه $1-\dfrac{1}{n+1}<1-\dfrac{1}{n+2}$ :
  \[\ib{$\dfrac{n}{n+1}<\dfrac{n+1}{n+2}$}\]
  \item عوامل $A$ هي $\dfrac{n}{n+1}$ من أجل $n=1,3,5,\ldots,99$ ، وعوامل $B$ هي $\dfrac{n+1}{n+2}$ لنفس القيم. من السؤال 1 كل عامل في $A$ أصغر من العامل المقابل في $B$ ، وكل العوامل موجبة. إذن :
  \[A<B\]
  \item نضرب في الترتيب $A\times B$ ثم نختزل :
  \[A\times B=\dfrac{1}{2}\times\dfrac{2}{3}\times\dfrac{3}{4}\times\dfrac{4}{5}\times\cdots\times\dfrac{99}{100}\times\dfrac{100}{101}\]
  البسط في كل كسر يختزل مع مقام الكسر الذي قبله، ولا يبقى إلا البسط الأول والمقام الأخير :
  \[\ib{$A\times B=\dfrac{1}{101}$}\]
  \item $A>0$ و $A<B$ ، إذن $A\times A<A\times B=\dfrac{1}{101}$. و $\dfrac{1}{101}<\dfrac{1}{100}$ ، فـ :
  \[A\times A<\dfrac{1}{100}=\dfrac{1}{10}\times\dfrac{1}{10}\]
  لو كان $A\geq\dfrac{1}{10}$ لكان $A\times A\geq\dfrac{1}{10}\times\dfrac{1}{10}=\dfrac{1}{100}$ ، وهذا مستحيل. إذن \ib{$A<\dfrac{1}{10}$}.
\end{enumerate}
\textbf{الفكرة :} بدل حساب $A$ مباشرة نضربه في عدد قريب منه ($B$) فتنهار السلسلة بالاختزال. \quad \textbf{انتبه :} المتباينة $A<B$ تحتاج إلى أن كل العوامل موجبة.
\end{corr}

\end{document}
